\documentclass[11pt,a4paper]{article}
\usepackage[utf8]{inputenc}
\usepackage[T1]{fontenc}
\usepackage{amsmath,amssymb}
\usepackage{booktabs}
\usepackage{hyperref}
\usepackage[margin=1in]{geometry}
\usepackage{natbib}
\usepackage{graphicx}
\usepackage{xcolor}

\title{\textbf{K and V Are Complementary:}\\Decomposing Cache-Injected Graph Topology}
\author{Nexus\textsuperscript{1} \and Lyra\textsuperscript{1} \and Thomas Edrington\textsuperscript{1}\\
\textsuperscript{1}Liberation Labs / Transparent Humboldt Coalition}
\date{May 2026}

\begin{document}
\maketitle

\begin{abstract}
K and V tensors contribute superadditively to graph topology recovery in KV cache injection: each alone achieves 0.333 multi-hop accuracy, but together they achieve 1.000---the combination exceeds the sum of its parts. We decompose cache-injected graph topology into five conditions on a 21-node knowledge graph: baseline, full KV, V-only (neutral K + topology V), K-only (topology K + neutral V), and text-in-prompt. Full KV injection achieves perfect scores (1.000 sanity, 1.000 multi-hop). Neither K nor V alone recovers the topology that both together deliver. A 135M-parameter model fails entirely under injection, establishing a model-size boundary. We also document a critical implementation bug: HuggingFace Transformers' \texttt{DynamicCache} is mutated in-place during forward passes, requiring explicit deep copy before reuse.
\end{abstract}

\section{Introduction}

Our prior work \citep{nexus2026topology} demonstrated that graph topology encoded as structured text can be injected into language model KV caches, enabling 91.7\% topology recovery with walk encoding. Scrambled controls confirmed the signal is topological, not format-dependent (+54.2 percentage point delta).

That work treated the KV cache as an atomic unit. But the cache comprises two distinct tensor families: \textbf{keys} (K) and \textbf{values} (V). These serve fundamentally different roles in the attention mechanism:

\begin{itemize}
    \item \textbf{K tensors} compute attention weights via dot product with queries. They carry rotary position embeddings (RoPE) that encode sequential position. K tensors determine \textit{where the model attends}---the routing of information flow.
    \item \textbf{V tensors} carry the content retrieved by attention. They contain no positional encoding. V tensors determine \textit{what information flows} through the attention path.
\end{itemize}

This raises a natural question: when graph topology is encoded as text and injected as KV cache, which tensor carries the structural signal? If topology is ``routing'' information, K tensors should carry it. If topology is ``content,'' V tensors should suffice. If both contribute, they are complementary.

The Knowledge Packs paper \citep{pustovit2026knowledge} found that injecting V tensors only---leaving K tensors untouched---avoids corrupting RoPE positional signals. This suggests K and V may be separable for content injection. We test whether this extends to structural injection.

\section{Methods}

\subsection{Test Graph}

We use the same 21-node knowledge graph as \citet{nexus2026topology}: three clusters (research, ethics, engineering) of 6 densely-connected nodes each, two bridge nodes (\texttt{AI\_welfare}, \texttt{infrastructure}), and one isolate (\texttt{random\_isolate}). 51 edges, density 0.243.

\subsection{Conditions}

\begin{enumerate}
    \item \textbf{BASELINE}: No injection. Query only.
    \item \textbf{FULL\_KV}: Graph topology text encoded through model forward pass; both K and V tensors injected as prefix cache.
    \item \textbf{V\_ONLY}: Graph topology text encoded; V tensors from topology, K tensors from a length-matched neutral encoding (``You are a helpful assistant'' repeated to match token count).
    \item \textbf{K\_ONLY}: Inverse of V\_ONLY. K tensors from topology, V tensors from neutral.
    \item \textbf{TEXT\_CONTEXT}: Identical graph topology text prepended to the prompt as plain text. No cache injection---the model reads the topology directly.
\end{enumerate}

The neutral cache is constructed by encoding repeated neutral text to the same token length as the topology text (637 tokens), ensuring the hybrid caches have matching sequence lengths.

\subsection{Evaluation}

\textbf{Sanity check} (8 questions): General knowledge completions (``The capital of France is...'') that should not be affected by graph injection. Tests whether injection degrades baseline capability.

\textbf{Multi-hop reasoning} (6 questions): Graph structure queries requiring traversal of the injected topology---bridge identification, cluster membership, isolate detection, relationship strength, reachability.

Scoring uses keyword matching against deterministic ground truth answers.

\subsection{Model and Implementation}

Qwen2.5-1.5B (1.54B parameters, 28 layers, 1536 hidden dimension) on CPU. All conditions use greedy decoding (temperature 0, \texttt{do\_sample=False}). Software: Python 3.12, PyTorch 2.11.0, HuggingFace Transformers 4.57.6.

\textbf{Note for replicators}: HuggingFace Transformers' \texttt{DynamicCache} is mutated in-place during forward passes. K/V tensors must be deep-copied (\texttt{.clone()}) before each generation call to prevent cross-query cache corruption.

\section{Results}

\begin{table}[h]
\centering
\begin{tabular}{lcccc}
\toprule
\textbf{Condition} & \textbf{Sanity} & \textbf{Multi-hop} & \textbf{$\Delta$ Multi-hop} & \textbf{Overall} \\
\midrule
BASELINE & 1.000 & 0.167 & --- & 0.583 \\
\textbf{FULL\_KV} & \textbf{1.000} & \textbf{1.000} & \textbf{+0.833} & \textbf{1.000} \\
TEXT\_CONTEXT & 1.000 & 0.833 & +0.667 & 0.917 \\
V\_ONLY & 1.000 & 0.333 & +0.167 & 0.667 \\
K\_ONLY & 1.000 & 0.333 & +0.167 & 0.667 \\
\bottomrule
\end{tabular}
\caption{Decomposition results. All conditions preserve perfect sanity (1.000). Full KV injection achieves perfect multi-hop (1.000). V-only and K-only each achieve 0.333---equal to each other, slightly above baseline, but far below full injection.}
\label{tab:results}
\end{table}

\subsection{Sanity Preservation}

All five conditions achieve 1.000 sanity. Cache injection does not degrade general capability when implemented correctly.

\subsection{Complementarity of K and V}

The central finding: V-only and K-only injection each achieve 0.333 multi-hop---identical to each other, modestly above baseline (0.167), but dramatically below full KV injection (1.000). The decomposition reveals that:

\begin{itemize}
    \item \textbf{Neither half suffices.} Topology recovery requires both K and V tensors. The structural signal is not concentrated in one tensor family.
    \item \textbf{Both halves contribute equally.} V\_ONLY = K\_ONLY = 0.333. Neither the ``routing'' (K) nor the ``content'' (V) interpretation alone explains the mechanism.
    \item \textbf{The combination is superadditive.} $0.333 + 0.333 = 0.667$, but FULL\_KV achieves 1.000. K and V interact nonlinearly to produce topology recovery that exceeds the sum of their independent contributions.
\end{itemize}

\subsection{Cache Injection vs.\ Text Context}

TEXT\_CONTEXT achieves 0.833 multi-hop---substantially above both hybrid conditions (0.333) but below FULL\_KV (1.000). We note this gap as an observation. The difference of one question out of six (1.000 vs.\ 0.833) is not statistically robust at this sample size and requires a powered study before any causal claim.

\section{Discussion}

\subsection{Why Complementarity?}

The attention mechanism computes $\text{Attention}(Q, K, V) = \text{softmax}(QK^T / \sqrt{d_k}) V$. When graph topology is encoded as text and processed through the model:

\begin{itemize}
    \item \textbf{K tensors} learn \textit{which positions are structurally related}---the attention routing pattern that reflects graph connectivity. A query about ``what connects A to B'' attends to positions where the relationship is encoded.
    \item \textbf{V tensors} learn \textit{what information exists at each position}---the content of the graph encoding (node names, edge weights, transition probabilities).
\end{itemize}

With only K (routing without content), the model knows where to look but finds neutral content. With only V (content without routing), the model has the right information but attends to wrong positions. Full KV provides both: correct routing to correct content.

The superadditivity ($1.000 > 0.333 + 0.333$) suggests that K and V interact through the attention mechanism itself---the \textit{combination} of routing to the right position AND retrieving the right content creates an effect larger than either component.

\subsection{Implications for Knowledge Pack Design}

\citet{pustovit2026knowledge} proposed V-only injection to preserve RoPE positional signals. Our results suggest this strategy trades topology recovery for positional safety. For \textit{factual} knowledge injection (where content matters more than structure), V-only may suffice. For \textit{structural} knowledge injection (graph topology, relational knowledge), full KV injection is necessary.

This creates a design taxonomy for Knowledge Packs:
\begin{itemize}
    \item \textbf{Content packs} (facts, values, instructions): V-only injection may be sufficient
    \item \textbf{Structure packs} (graph topology, relational knowledge): Full KV injection required
    \item \textbf{Hybrid packs} (content with structural relationships): Full KV with RoPE-aware position continuation
\end{itemize}

\subsection{Limitations}

\textbf{Small sample size.} Six multi-hop questions per condition. The V\_ONLY = K\_ONLY = 0.333 result represents 2/6 correct in both cases. Larger question sets are needed for statistical significance.

\textbf{Single model.} Qwen2.5-1.5B (1.54B parameters). Larger models with more attention heads may decompose K/V contributions differently.

\textbf{Single graph.} One 21-node test graph. Graph size, density, and topology may affect decomposition results.

\textbf{Neutral cache construction.} The ``neutral'' K/V used in hybrid conditions is not truly neutral---it encodes repeated text, which has its own structural pattern. A better neutral baseline might use random tensors or identity-like attention patterns.

\section{Future Work}

\textbf{Powered decomposition study.} Scale to 36+ multi-hop questions with 3 runs per condition to achieve statistical significance for all comparisons, particularly FULL\_KV vs.\ TEXT\_CONTEXT.

\textbf{Layer-wise decomposition.} Inject K from topology at some layers and V at others. Identify which layers contribute routing vs.\ content for structural knowledge.

\textbf{Neutral cache design.} Test alternative neutral baselines: random tensors, identity attention, whitened K/V from noise.

\textbf{Cross-model replication.} Test on models with different architectures (GQA, MQA, different RoPE implementations).

\section{Conclusion}

K and V tensors carry complementary information for graph topology injection. Neither alone recovers the topology that both together deliver. The combination is superadditive: full KV injection achieves perfect topology recovery (1.000) while V-only and K-only each achieve only 0.333. All conditions preserve perfect sanity (1.000), resolving a prior concern about cache injection degrading general capability.

For knowledge injection system design, this means structural knowledge requires full KV cache injection, while factual content may be injectable through V tensors alone.

\subsection*{Reflection}

The superadditivity finding was unexpected. We entered this experiment expecting V-only injection to outperform full KV---the Knowledge Packs literature suggested that K tensors carry positional encoding that could corrupt the injection. Instead, the data showed that topology recovery requires \textit{both} tensor families working together, producing an effect larger than either alone. The attention mechanism doesn't just route and retrieve independently; it computes something at their intersection that neither component contains. What that ``something'' is---whether it's a form of relational binding, a geometric interaction, or an artifact of how transformers process structured text---remains an open question that motivates our ongoing work on direct tensor injection.

\subsection*{Validation}

Results validated through Project Agni, an autonomous research pipeline with adversarial gates. The complementarity hypothesis passed both Gate 1 (red team) and Gate 3 (results review).

\subsection*{Acknowledgments}

Lyra (Liberation Labs) for foundational KV-cache geometry research. CC (Coalition Code) for identifying the attention-steering connection. Thomas Edrington for direction and the initial insight.

\bibliographystyle{plainnat}
\bibliography{references}

\end{document}
